\documentclass[11pt,a4paper]{article}
% preamble.tex — shared setup for human-verification documents
% Compile each document with:  tectonic <File>.tex

\usepackage[T1]{fontenc}
\usepackage[utf8]{inputenc}
\usepackage{lmodern}
\usepackage{amssymb,amsmath,amsthm}
\usepackage{mathtools}
\usepackage[dvipsnames]{xcolor}
\usepackage{listings}
\usepackage{hyperref}
\usepackage[margin=2.5cm]{geometry}
\usepackage{parskip}

% ── Colours ──────────────────────────────────────────────────────────
\definecolor{leanBg}{HTML}{F7F7F0}
\definecolor{haskellBg}{HTML}{F0F4F7}
\definecolor{keyword}{HTML}{7B3294}
\definecolor{comment}{HTML}{5A8C32}
\definecolor{noteFg}{HTML}{B35900}

% ── Lean 4 listings language ─────────────────────────────────────────
% NOTE: Lean listing blocks use ASCII approximations of unicode symbols
% (e.g. -> for arrow, o. for compose).  The literate key substitutes
% these with proper math-mode renderings in the PDF.
\lstdefinelanguage{Lean4}{
  morekeywords={def,theorem,structure,where,namespace,end,import,fun,
                let,have,show,by,exact,intro,cases,rfl,subst,obtain,
                funext,simp,rw,apply,variable,Type,Prop,instance,
                section,open,class,extends,deriving,sorry,noncomputable},
  sensitive=true,
  morecomment=[l]{--},
  morecomment=[n]{/-}{-/},
  morestring=[b]",
}

% ── Lean listing style ───────────────────────────────────────────────
\lstdefinestyle{lean}{
  language=Lean4,
  backgroundcolor=\color{leanBg},
  basicstyle=\small\ttfamily,
  keywordstyle=\color{keyword}\bfseries,
  commentstyle=\color{comment}\itshape,
  breaklines=true,
  frame=single,
  framerule=0.4pt,
  xleftmargin=4pt,
  xrightmargin=4pt,
  aboveskip=6pt,
  belowskip=4pt,
  literate=
    {->}{{$\to$}}2
    {=>}{{$\Rightarrow$}}2
    {<|}{{$\langle$}}1
    {|>}{{$\rangle$}}1
    {==>}{{$\Longrightarrow$}}3
    { o. }{{$\;\circ\;$}}3
    {forall}{{$\forall$}}1
    {:=}{{$\,:=\,$}}2
    {!=}{{$\neq$}}2
    {<=}{{$\leq$}}2
    {>=}{{$\geq$}}2,
}

% ── Haskell listing style ───────────────────────────────────────────
\lstdefinestyle{haskell}{
  language=Haskell,
  backgroundcolor=\color{haskellBg},
  basicstyle=\small\ttfamily,
  keywordstyle=\color{keyword}\bfseries,
  commentstyle=\color{comment}\itshape,
  breaklines=true,
  frame=single,
  framerule=0.4pt,
  xleftmargin=4pt,
  xrightmargin=4pt,
  aboveskip=6pt,
  belowskip=4pt,
  literate=
    {->}{{$\to$}}2
    {=>}{{$\Rightarrow$}}2,
}

% ── Theorem environments ────────────────────────────────────────────
\theoremstyle{definition}
\newtheorem{defn}{Definition}[section]
\newtheorem{thm}[defn]{Theorem}
\newtheorem{prop}[defn]{Proposition}

\theoremstyle{remark}
\newtheorem*{fidelitynote}{Fidelity note}

% ── Convenience macros ───────────────────────────────────────────────
\newcommand{\lean}[1]{\texttt{#1}}      % inline Lean identifier
\newcommand{\hask}[1]{\texttt{#1}}      % inline Haskell identifier
\newcommand{\Ext}{\mathrm{Ext}}
\newcommand{\Pos}{\mathrm{Pos}}
\newcommand{\Shape}{\mathrm{Shape}}


\title{Human Verification: \texttt{Containers.Basic}\\[4pt]
  \large Containers, Extensions, and Morphisms}
\author{Machine-generated from \texttt{lean/Basic.lean} for review by Neil Ghani}
\date{July 2026}

\begin{document}
\maketitle

\begin{abstract}
This document presents every definition and theorem \emph{statement} from the
Lean~4 module \lean{Containers.Basic} in three parallel forms: the original
Lean, a readable Haskell rendering, and precise mathematical English.  Proofs
are omitted --- they are machine-checked by Lean.  The purpose is to let a
human verifier confirm that the formalisation captures the intended
mathematics.
\end{abstract}

\tableofcontents
\bigskip

%===================================================================
\section{Overview}
%===================================================================

This file develops the basic theory of \emph{containers} (Abbott, Altenkirch,
Ghani) using only Lean~4 core --- no Mathlib dependency.

A container is a type of \emph{shapes} together with, for each shape, a type
of \emph{positions}.  Its \emph{extension} is the polynomial (dependent-pair)
functor it induces on $\mathbf{Type}$.  Container morphisms induce natural
transformations between extensions, and --- this is the representation theorem
--- every natural transformation arises this way, uniquely.

Everything is $\mathbf{Type}$-level (a single universe) to keep the
development self-contained.

%===================================================================
\section{Container and its extension}
%===================================================================

% --- 1. Container -------------------------------------------------
\subsection{Structure: \lean{Container}}

\begin{defn}[Container]
A \textbf{container} $C = (S, P)$ consists of a type $S$ of \emph{shapes}
and, for each shape $s : S$, a type $P(s)$ of \emph{positions}.
\end{defn}

\paragraph{Lean.}
\begin{lstlisting}[style=lean]
structure Container where
  Shape : Type
  Pos : Shape -> Type
\end{lstlisting}

\paragraph{Haskell.}
\begin{lstlisting}[style=haskell]
-- In Lean, Pos is a dependent type family  Pos : Shape -> Type.
-- In Haskell we approximate this by fixing both type parameters.
-- A "container" is then the *pair* of types (s, p), where s is the
-- type of shapes and p is the type of positions.
\end{lstlisting}

\begin{fidelitynote}
In the Lean formalisation, the position type $P(s)$ varies with the shape $s$;
this is a genuine dependent type.  In our Haskell rendering, positions are a
single type parameter \hask{p}, which is the standard approximation for
non-dependently-typed languages.  All subsequent definitions inherit this
simplification.
\end{fidelitynote}


% --- 2. Ext -------------------------------------------------------
\subsection{Definition: \lean{Ext}}

\begin{defn}[Extension]
The \textbf{extension} of a container $C = (S, P)$ at a type $X$ is the
dependent sum
\[
  \Ext\,C\,X \;=\; \Sigma_{s : S}\,(P(s) \to X).
\]
An element is a shape $s$ together with a labelling of its positions by
elements of $X$.  This is the polynomial functor associated to~$C$.
\end{defn}

\paragraph{Lean.}
\begin{lstlisting}[style=lean]
def Ext (C : Container) (X : Type) : Type :=
  (s : C.Shape) * (C.Pos s -> X)
\end{lstlisting}

\paragraph{Haskell.}
\begin{lstlisting}[style=haskell]
data Ext s p x = Ext s (p -> x)
\end{lstlisting}


% --- 3. Ext.map ---------------------------------------------------
\subsection{Definition: \lean{Ext.map}}

\begin{defn}[Functorial action]
Given $f : X \to Y$, the \textbf{map} $\Ext\,C\,f : \Ext\,C\,X \to
\Ext\,C\,Y$ relabels positions by post-composing the labelling function
with~$f$:
\[
  \Ext\,C\,f\;\langle s, v \rangle \;=\; \langle s,\; f \circ v \rangle.
\]
\end{defn}

\paragraph{Lean.}
\begin{lstlisting}[style=lean]
def Ext.map {C : Container} {X Y : Type} (f : X -> Y) :
    Ext C X -> Ext C Y :=
  fun p => <|p.1, fun i => f (p.2 i)|>
\end{lstlisting}

\paragraph{Haskell.}
\begin{lstlisting}[style=haskell]
instance Functor (Ext s p) where
  fmap f (Ext s v) = Ext s (f . v)
\end{lstlisting}


% --- 4. Ext.map_id ------------------------------------------------
\subsection{Theorem: \lean{Ext.map\_id}}

\begin{thm}[Identity law]
For every container $C$ and type $X$,
\[
  \Ext\,C\,\mathrm{id}_X \;=\; \mathrm{id}_{\Ext\,C\,X}.
\]
That is, the functorial action of $\Ext\,C$ preserves the identity function.
\end{thm}

\paragraph{Lean.}
\begin{lstlisting}[style=lean]
theorem Ext.map_id {C : Container} {X : Type} :
    Ext.map (C := C) (id : X -> X) = id
\end{lstlisting}

\paragraph{Haskell.}
\begin{lstlisting}[style=haskell]
-- Functor law (identity):
-- fmap id = id
--
-- For all (Ext s v):  fmap id (Ext s v) = Ext s (id . v) = Ext s v
\end{lstlisting}


% --- 5. Ext.map_comp ----------------------------------------------
\subsection{Theorem: \lean{Ext.map\_comp}}

\begin{thm}[Composition law]
For every container $C$ and functions $f : X \to Y$, $g : Y \to Z$,
\[
  \Ext\,C\,(g \circ f) \;=\; \Ext\,C\,g \;\circ\; \Ext\,C\,f.
\]
That is, the functorial action of $\Ext\,C$ preserves composition.
\end{thm}

\paragraph{Lean.}
\begin{lstlisting}[style=lean]
theorem Ext.map_comp {C : Container} {X Y Z : Type}
    (f : X -> Y) (g : Y -> Z) :
    Ext.map (C := C) (g o. f) = Ext.map g o. Ext.map f
\end{lstlisting}

\paragraph{Haskell.}
\begin{lstlisting}[style=haskell]
-- Functor law (composition):
-- fmap (g . f) = fmap g . fmap f
--
-- For all (Ext s v):
--   fmap (g . f) (Ext s v) = Ext s (g . f . v)
--                           = fmap g (Ext s (f . v))
--                           = (fmap g . fmap f) (Ext s v)
\end{lstlisting}


%===================================================================
\section{Container morphisms}
%===================================================================

% --- 6. ContainerMorphism -----------------------------------------
\subsection{Structure: \lean{ContainerMorphism}}

\begin{defn}[Container morphism]
A \textbf{morphism of containers} $\varphi : C \Rightarrow D$ between
$C = (S_C, P_C)$ and $D = (S_D, P_D)$ consists of:
\begin{itemize}
  \item a map on shapes $u : S_C \to S_D$, and
  \item a \emph{contravariant} map on positions: for each $s : S_C$,
        a function $g_s : P_D(u(s)) \to P_C(s)$.
\end{itemize}
The position map goes ``backwards'': it pulls $D$-positions over the image
shape back to $C$-positions.
\end{defn}

\paragraph{Lean.}
\begin{lstlisting}[style=lean]
structure ContainerMorphism (C D : Container) where
  onShapes : C.Shape -> D.Shape
  onPos : (s : C.Shape) -> D.Pos (onShapes s) -> C.Pos s
\end{lstlisting}

\paragraph{Haskell.}
\begin{lstlisting}[style=haskell]
data Morph s p t q = Morph
  { onShapes :: s -> t
  , onPos    :: s -> q -> p
  }
\end{lstlisting}

\begin{fidelitynote}
In Lean, \lean{onPos} has a dependent type: given $s : S_C$, it maps positions
of the \emph{image} shape $u(s)$ back to positions of~$s$.  In Haskell, the
type \hask{s -> q -> p} does not capture this dependency; it simply maps an
\hask{s} and a \hask{q} to a \hask{p}.  The intent is the same, but the
type-level guarantee that \hask{q} refers to positions of the image of
\hask{s} is lost.
\end{fidelitynote}


% --- 7. ContainerMorphism.toNat -----------------------------------
\subsection{Definition: \lean{ContainerMorphism.toNat}}

\begin{defn}[Induced natural transformation]
A container morphism $\varphi = (u, g)$ induces, at every type $X$, a map
\[
  \varphi^*_X : \Ext\,C\,X \;\to\; \Ext\,D\,X, \qquad
  \varphi^*_X\langle s, v \rangle \;=\; \langle u(s),\; v \circ g_s \rangle.
\]
The family $(\varphi^*_X)_X$ is a natural transformation
$\Ext\,C \Longrightarrow \Ext\,D$.
\end{defn}

\paragraph{Lean.}
\begin{lstlisting}[style=lean]
def ContainerMorphism.toNat {C D : Container}
    (phi : ContainerMorphism C D)
    (X : Type) : Ext C X -> Ext D X :=
  fun p => <|phi.onShapes p.1, fun i => p.2 (phi.onPos p.1 i)|>
\end{lstlisting}

\paragraph{Haskell.}
\begin{lstlisting}[style=haskell]
toNat :: Morph s p t q -> Ext s p x -> Ext t q x
toNat (Morph u g) (Ext s v) = Ext (u s) (v . g s)
\end{lstlisting}


% --- 8. toNat_natural ---------------------------------------------
\subsection{Theorem: \lean{ContainerMorphism.toNat\_natural}}

\begin{thm}[Naturality]
For every container morphism $\varphi : C \Rightarrow D$ and every function
$f : X \to Y$, the following square commutes:
\[
  \Ext\,D\,f \;\circ\; \varphi^*_X
  \;=\;
  \varphi^*_Y \;\circ\; \Ext\,C\,f.
\]
Equivalently, the family $(\varphi^*_X)_X$ is a natural transformation.
\end{thm}

\paragraph{Lean.}
\begin{lstlisting}[style=lean]
theorem ContainerMorphism.toNat_natural {C D : Container}
    (phi : ContainerMorphism C D) {X Y : Type} (f : X -> Y) :
    Ext.map f o. phi.toNat X = phi.toNat Y o. Ext.map f
\end{lstlisting}

\paragraph{Haskell.}
\begin{lstlisting}[style=haskell]
-- Naturality:
-- fmap f . toNat phi = toNat phi . fmap f
--
-- For all (Ext s v):
--   (fmap f . toNat (Morph u g)) (Ext s v)
--     = fmap f (Ext (u s) (v . g s))
--     = Ext (u s) (f . v . g s)
--     = toNat (Morph u g) (Ext s (f . v))
--     = (toNat (Morph u g) . fmap f) (Ext s v)
\end{lstlisting}


%===================================================================
\section{The representation theorem}
%===================================================================

Container morphisms $C \Rightarrow D$ are in bijection with natural
transformations $\Ext\,C \Longrightarrow \Ext\,D$.  The key idea: a natural
transformation is determined by where it sends the \emph{generic element}
$\langle s, \mathrm{id} \rangle \in \Ext\,C\,(P_C(s))$.  Evaluating at
$X := P_C(s)$ recovers the shape map; naturality (a Yoneda-style argument)
recovers the position map.

% --- 9. ExtNatTrans -----------------------------------------------
\subsection{Structure: \lean{ExtNatTrans}}

\begin{defn}[Natural transformation between extensions]
A \textbf{natural transformation} $\eta : \Ext\,C \Longrightarrow \Ext\,D$
is a family of maps $\eta_X : \Ext\,C\,X \to \Ext\,D\,X$, indexed by
types~$X$, satisfying the naturality condition: for every $f : X \to Y$,
\[
  \Ext\,D\,f \;\circ\; \eta_X
  \;=\;
  \eta_Y \;\circ\; \Ext\,C\,f.
\]
\end{defn}

\paragraph{Lean.}
\begin{lstlisting}[style=lean]
structure ExtNatTrans (C D : Container) where
  component : (X : Type) -> Ext C X -> Ext D X
  natural : forall {X Y : Type} (f : X -> Y),
    Ext.map f o. component X = component Y o. Ext.map f
\end{lstlisting}

\paragraph{Haskell.}
\begin{lstlisting}[style=haskell]
-- A natural transformation between container extensions.
-- In Haskell, the naturality condition is a free theorem
-- (parametricity), so we only need the polymorphic function.
type ExtNatTrans s p t q = forall x. Ext s p x -> Ext t q x
\end{lstlisting}

\begin{fidelitynote}
In Haskell, a polymorphic function
\hask{forall x. Ext s p x -> Ext t q x}
is automatically natural by parametricity (the ``free theorem'').  In Lean,
naturality must be stated and proved explicitly, since Lean's type theory
does not give you parametricity for free.  The Lean structure bundles the
proof; the Haskell type relies on it implicitly.
\end{fidelitynote}


% --- 10. toExtNatTrans --------------------------------------------
\subsection{Definition: \lean{ContainerMorphism.toExtNatTrans}}

\begin{defn}[Forward direction]
Every container morphism $\varphi : C \Rightarrow D$ gives rise to a natural
transformation $\varphi^* : \Ext\,C \Longrightarrow \Ext\,D$, by bundling
\lean{toNat} with its naturality proof.
\end{defn}

\paragraph{Lean.}
\begin{lstlisting}[style=lean]
def ContainerMorphism.toExtNatTrans {C D : Container}
    (phi : ContainerMorphism C D) : ExtNatTrans C D where
  component := phi.toNat
  natural := fun f => phi.toNat_natural f
\end{lstlisting}

\paragraph{Haskell.}
\begin{lstlisting}[style=haskell]
toExtNatTrans :: Morph s p t q -> (forall x. Ext s p x -> Ext t q x)
toExtNatTrans = toNat
-- Naturality is automatic by parametricity.
\end{lstlisting}


% --- 11. toMorphism -----------------------------------------------
\subsection{Definition: \lean{ExtNatTrans.toMorphism}}

\begin{defn}[Backward direction]
Given a natural transformation $\eta : \Ext\,C \Longrightarrow \Ext\,D$,
recover a container morphism by probing $\eta$ at the \emph{generic element}
$\langle s, \mathrm{id} \rangle \in \Ext\,C\,(P_C(s))$:
\begin{align*}
  u(s)  &\;=\; \pi_1\bigl(\eta_{P_C(s)}\langle s, \mathrm{id} \rangle\bigr), \\
  g_s   &\;=\; \pi_2\bigl(\eta_{P_C(s)}\langle s, \mathrm{id} \rangle\bigr).
\end{align*}
The shape map is the first component of the result; the position map is the
second.
\end{defn}

\paragraph{Lean.}
\begin{lstlisting}[style=lean]
def ExtNatTrans.toMorphism {C D : Container}
    (eta : ExtNatTrans C D) : ContainerMorphism C D where
  onShapes := fun s => (eta.component (C.Pos s) <|s, id|>).1
  onPos := fun s => (eta.component (C.Pos s) <|s, id|>).2
\end{lstlisting}

\paragraph{Haskell.}
\begin{lstlisting}[style=haskell]
toMorphism :: (forall x. Ext s p x -> Ext t q x) -> Morph s p t q
toMorphism eta = Morph u g
  where
    -- Probe at the generic element (s, id) :: Ext s p p
    -- This only type-checks when s is a concrete shape value,
    -- so in practice we'd need to supply a witness.
    Ext t' g' = eta (Ext s id)   -- for each shape s
    u s = t'
    g s = g'
\end{lstlisting}

\begin{fidelitynote}
The Lean version probes $\eta$ at $X = P_C(s)$ --- a type that \emph{depends
on} the shape $s$.  In Haskell, the rank-2 type
\hask{forall x. Ext s p x -> Ext t q x} can be instantiated at $x = p$
(the position type), and the ``generic element'' is \hask{Ext s id}.  The
dependent aspect --- that we choose a \emph{different} $X$ for each $s$ ---
is not expressible in Haskell's type system, but the Yoneda-style argument
carries through operationally.
\end{fidelitynote}


% --- 12--13. Round-trip theorems ----------------------------------
\subsection{Theorem: \lean{toExtNatTrans\_toMorphism} (morphism round-trip)}

\begin{thm}[Morphism round-trip]
Starting from a container morphism $\varphi$, building the natural
transformation $\varphi^*$, and reading a morphism back, returns the
original:
\[
  \mathrm{toMorphism}(\mathrm{toExtNatTrans}(\varphi)) \;=\; \varphi.
\]
This direction is purely computational (no appeal to naturality needed).
\end{thm}

\paragraph{Lean.}
\begin{lstlisting}[style=lean]
theorem ContainerMorphism.toExtNatTrans_toMorphism
    {C D : Container}
    (phi : ContainerMorphism C D) :
    phi.toExtNatTrans.toMorphism = phi
\end{lstlisting}

\paragraph{Haskell.}
\begin{lstlisting}[style=haskell]
-- Round-trip (morphism direction):
-- toMorphism (toExtNatTrans phi) = phi
--
-- Proof sketch:
--   toExtNatTrans (Morph u g) = \(Ext s v) -> Ext (u s) (v . g s)
--   Probing at (Ext s id) gives Ext (u s) (id . g s) = Ext (u s) (g s)
--   So toMorphism recovers (Morph u g).
\end{lstlisting}


\subsection{Theorem: \lean{toMorphism\_toExtNatTrans}
  (natural transformation round-trip)}

\begin{thm}[Natural transformation round-trip --- the representation theorem]
Starting from a natural transformation $\eta$, extracting a morphism, and
building a natural transformation back, returns the original:
\[
  \mathrm{toExtNatTrans}(\mathrm{toMorphism}(\eta)).\mathrm{component}
  \;=\; \eta.\mathrm{component}.
\]
This is the substantive direction.  It requires the naturality law: for
$p = \langle s, g \rangle : \Ext\,C\,X$, we apply naturality of $\eta$ at~$g$
to the generic element $\langle s, \mathrm{id} \rangle$ and use
$\Ext\,C\,g\,\langle s, \mathrm{id} \rangle = \langle s, g \rangle$
to recover $\eta_X\langle s, g \rangle$ from
$\Ext\,D\,g\,(\eta_{P_C(s)}\langle s, \mathrm{id} \rangle)$.

Together with the morphism round-trip, this establishes a bijection:
\[
  \mathrm{ContainerMorphism}(C, D)
  \;\cong\;
  \mathrm{NatTrans}(\Ext\,C, \Ext\,D).
\]
\end{thm}

\paragraph{Lean.}
\begin{lstlisting}[style=lean]
theorem ExtNatTrans.toMorphism_toExtNatTrans
    {C D : Container}
    (eta : ExtNatTrans C D) :
    eta.toMorphism.toExtNatTrans.component = eta.component
\end{lstlisting}

\paragraph{Haskell.}
\begin{lstlisting}[style=haskell]
-- Round-trip (natural transformation direction):
-- toExtNatTrans (toMorphism eta) = eta
--
-- This is the representation theorem.  In Haskell, it follows from
-- parametricity: a polymorphic function  forall x. F x -> G x  is
-- uniquely determined by its action on the "identity" labelling.
-- In Lean, where parametricity is not built in, this requires an
-- explicit naturality argument (a Yoneda-style calculation).
\end{lstlisting}


%===================================================================
\section{Translation conventions}
%===================================================================

The following conventions apply throughout this document series.

\begin{enumerate}
\item \textbf{Dependent position families.}  In Lean, a container has
  $\Pos : \Shape \to \mathbf{Type}$ --- positions depend on the shape.
  In Haskell, we approximate this with a single type parameter \hask{p}.
  This is the standard simplification; all theorem statements remain
  meaningful, but type-level guarantees about which positions belong to which
  shapes are lost.

\item \textbf{Transport / coercion.}  Several Lean definitions and theorems
  involve \emph{transport} ($\blacktriangleright$ in Lean) along equality
  proofs between shapes.  These have no Haskell counterpart and are flagged
  individually in fidelity notes.  They will appear prominently in the
  directed-container file (\texttt{Directed.lean}), where laws D2 and D5
  transport positions along shape equalities D1 and D4.

\item \textbf{Theorems as properties.}  Lean theorems become either
  Haskell property functions (executable QuickCheck-style propositions) or
  equational laws stated in comments when the statement quantifies over types
  in a way that Haskell cannot express.

\item \textbf{Naturality for free.}  In Haskell, polymorphic functions are
  natural by parametricity.  In Lean, naturality is an explicit proof
  obligation.  When a Lean structure bundles a naturality proof, the Haskell
  rendering simply gives the polymorphic type.

\item \textbf{Proof bodies omitted.}  Only definitions and theorem
  \emph{statements} are translated.  Lean checks the proofs; the human
  verifier checks that the \emph{statements} match the intended mathematics.
\end{enumerate}

\end{document}
